\section{Conclusion}
\label{sec:conclusion}

This project investigated how a compact visual theme can be converted into a
mapped, technically valid, and reproducible game-asset texture. The resulting
\skinsmith\ prototype combines three human checkpoints with asset-aware
generation, deterministic UV baking, multi-view evaluation, hard feasibility
rules, and one bounded refinement round.

For RQ1, the four-pair AK-47 experiment shows that composing in continuous
weapon space and baking into the existing atlas improves both true asset-seam and
multi-view performance relative to the tileable baseline in every observed pair.
For RQ2, mapped evaluation exposes view-specific loss, internal seam failure, and
the difference between an attractive source and a usable texture. For RQ3, the
selection-policy ablation and the Route C results show that hard constraints and
rollback prevent acceptance of infeasible or lower-scoring candidates. For RQ4,
the live run demonstrates a complete path from a short theme through theme,
direction, and mapped-artwork decisions to a 2048 PNG/TGA export. For RQ5, the
M4A4 case shows that the core pipeline transfers through an adapter, while UV
topology and correction parameters remain asset-specific.

The main conclusion is methodological: image generation should be treated as one
stage in an asset pipeline, not as the endpoint. A credible generative Agent must
show what happens after mapping, preserve the user's authority where metrics are
weak, distinguish feasibility from preference, and retain the evidence behind
acceptance and rollback.

The current system is deliberately limited to base-colour Custom Paint Job
textures. It does not generate normal, roughness, metallic, height, or displacement
channels; Workbench lighting must therefore not be interpreted as generated
surface detail. Before final submission, the highest-value evidence addition is a
fixed left/right/top Workbench capture of the exact accepted dragon TGA, together
with one screenshot of the import settings.

Future work should proceed in four stages. First, add spatially aligned normal,
roughness, and metallic generation with cross-channel and engine-side validation.
Second, automate fixed-camera Workbench or game-engine capture so that deployment
evidence is reproducible. Third, evaluate more weapon and non-weapon adapters and
learn asset-specific correction parameters. Fourth, improve perceptual measures
and conduct a user study of the three checkpoints, including task time, revision
count, perceived control, and disagreements between source-only and mapped
choices. These extensions would test generalisation without weakening the bounded
claims established here.
